\documentclass{umk-matematika}
\begin{document}

\begin{center}
{\Large\bfseries Практическая работа 11}
\end{center}
\vspace{0.5cm}

\textbf{Задача 1.} Сформулируйте признак параллельности прямой и плоскости.\par\vspace{10pt}
\textbf{Задача 2.} Сформулируйте признак параллельности двух плоскостей.\par\vspace{10pt}
\textbf{Задача 3.} Какие прямые в пространстве называются скрещивающимися?\par\vspace{10pt}
\textbf{Задача 4.} В кубе $ABCDA\_1B\_1C\_1D\_1$ укажите рёбра, параллельные ребру $AB$.\par\vspace{10pt}
\textbf{Задача 5.} Дана треугольная пирамида $SABC$. Точки M и N — середины рёбер SA и SB соответственно. Докажите, что прямая MN параллельна плоскости ABC.\par\vspace{10pt}
\textbf{Задача 6.} Даны две параллельные плоскости. Прямая пересекает первую плоскость. Пересечёт ли она вторую?\par\vspace{10pt}
\textbf{Задача 7.} Через концы отрезка AB и его середину M проведены параллельные прямые, пересекающие плоскость $\alpha$ в точках A$_1$, B$_1$, M$_1$. Найдите MM$_1$, если AA$_1$ = 5 см, BB$_1$ = 11 см, и отрезок AB не пересекает плоскость.\par\vspace{10pt}
\textbf{Задача 8.} В кубе $ABCDA\_1B\_1C\_1D\_1$ точка K — середина ребра AA$_1$, точка L — середина ребра CC$_1$. Докажите, что прямая KL параллельна плоскости ABCD.\par\vspace{10pt}
\textbf{Задача 9.} Даны две скрещивающиеся прямые a и b. Через точку M, не лежащую на них, проведите плоскость, параллельную обеим прямым. Всегда ли это возможно?\par\vspace{10pt}
\textbf{Задача 10.} Две балки перекрытия лежат в параллельных плоскостях. Расстояние между плоскостями — 3 м. Концы верхней балки проектируются на нижнюю плоскость. Найдите расстояние между проекциями концов верхней балки, если её длина 5 м, а угол между балками 60$^\circ$.\par\vspace{10pt}
\newpage
\section*{Ответы}

\begin{enumerate}
  \item Ответ: признак параллельности прямой и плоскости.
  \item Ответ: признак параллельности плоскостей.
  \item Ответ: прямые, не лежащие в одной плоскости.
  \item Ответ: DC, A$_1$B$_1$, D$_1$C$_1$.
  \item Ответ: $MN \parallel (ABC)$.
  \item Ответ: да, пересечёт.
  \item Ответ: $MM\_1 = 8$ см.
  \item Ответ: $KL \parallel (ABCD)$.
  \item Ответ: да, всегда возможно.
  \item Ответ: 5 м.
\end{enumerate}
\end{document}
